水文頻度解析における確率分布の選択は, 設計外力の推定値に直接影響する重要な手続きである.
日本では, 水文頻度解析の考え方, 用いられる確率分布, 母数推定法および適合度評価法について, 早くから体系的な整理が行われてきた{\cite{Takara1998}}.
とくに, 宝・高棹は, 確率分布モデルの適合度と確率水文量の変動性を併せて評価する枠組みを提示し, 標準最小二乗規準（以下, SLSC）を実務的な適合度評価指標として位置づけた{\cite{TakaraTakasao1988}}.
\rev{SLSC は, 確率紙上の視覚的一致性を数量化しようとする指標であり, 確率紙上の直線性を相関係数で評価する Filliben の正規性検定や, Kinnison による Gumbel 分布の相関係数適合度検定とも同系統の発想に立つものとみなせる{\cite{Filliben1975,Kinnison1989}}.}

その一方で, SLSC の統計的意味や運用の妥当性については, 既にいくつかの問題点が指摘されている.
葛葉は, SLSC 自体が確率変数であり, 固定値による判定ではなく, 標本サイズに応じて閾値を考えるべきことを論じた{\cite{Kuzuha2010}}.
藤部は, モンテカルロ計算により, SLSC が標本サイズに依存して小さくなること, 固定閾値 $0.04$ による判定が安定的でないことを示した.
また, 再現期待値のバイアスと推定幅の比較から, 3母数分布間の選択による実用上の改善が限定的であることを指摘した{\cite{Fujibe2011}}.
林らは, SLSC を検定統計量として扱う試みを示したが, SLSC の分布が標本サイズや想定分布に依存し, 3母数型モデルに対して一般的な取扱いが困難であることを報告している{\cite{HayashiEtAl2012}}.
また, 葛葉・水木は, 従来の手引きに基づく手続きにおいて, SLSC が標本サイズおよび確率分布に関して一律には比較しにくいことを整理し, その修正法や適用上の留意点を提案している{\cite{KuzuhaMizuki2021,KuzuhaMizuki2022}}.

以上の既往研究により, SLSC の固定閾値運用や標本サイズ依存性の問題はかなり明らかになってきた.
しかし, SLSC の計算において, 観測値と理論値をいかなる空間で比較しているのか, すなわち標準化の役割そのものを整理した議論は十分ではない.
線形標準化であれば, それは位置と尺度の調整に過ぎず, 観測値の空間における距離評価と本質的に同値である.
これに対し, 対数変換を含む標準化のように非線形標準化を用いる場合には, SLSC の値は標準化関数に依存した距離評価となり, 分布間で同一の距離尺度を共有しているとは限らない.
しかし, 実務上は, 標準変量の定義の違いが SLSC の比較可能性に与える影響は必ずしも明示されてこなかった.
また, 実務における分布選択は, 単純な最小 SLSC ではなく, 一定の閾値を満たす候補の中からジャックナイフ法により推定幅の小さい分布を選ぶ手順として運用されることが多い{\cite{Fujibe2011}}.
したがって, 現在の実務に即した形で SLSC 系規準の位置づけを再検討するには, SLSC の値そのものの挙動に加え, SLSC+ジャックナイフ法および情報量規準による採択結果との比較が必要である.

そこで本稿では, SLSC の比較空間を二つに分けて考える.
すなわち, 観測値と理論分位を実データの空間で比較する立場を, 本稿では便宜的に X 空間, 標準化後の確率紙上の空間で比較する立場を S 空間と呼ぶ.
本稿の関心は, これらを明確に区別した上で, とくに S 空間における非線形標準化が適合度評価にどのような偏りを導くかを理論と数値計算の両面から再検討することにある.

本稿では, グンベル分布 (Gumbel), 一般化極値分布 (GEV), 対数ピアソンIII型分布 (LP3), 平方根指数型最大値分布 (SqrtEt), 平行移動指数分布 (Exp), 3母数対数正規分布 (LN3) の6分布を対象として, 多数回のモンテカルロ計算を実施した.
\rev{これらの分布は, 日本の実務で用いられてきた水文統計ユーティリティー Ver. 1.5 で扱われる分布を踏まえて選定した{\cite{JICE2003}}.}
本稿の目的は, 既往研究で指摘されてきた問題を踏まえつつ, 次の三点を明確にすることである.
第一に, SLSC の標本サイズ依存性が6分布比較に拡張しても維持されるかを確認する.
第二に, SLSC の値が線形変換では不変である一方, 非線形標準化では標準化関数に依存した距離評価となることを示す.
第三に, SLSC+ジャックナイフ法と AIC による採択結果を比較し, 実際の分布選択手順における SLSC 系規準の位置づけを再確認する.

以下, 2章で方法, 3章で結果, 4章で考察, 5章で結論を示す.
